\bookchapter{Agent, environment and reward}\label{ch:rl1}

\section*{Explanation}
In most of this book so far a network learned from \emph{examples with the right answer}: an image together with the label ``ball''. For many robot tasks nobody can give the right answer for every moment. Nobody can tell a walking robot the exact angles of its twelve leg joints for every millisecond. What we \emph{can} do is judge the result: ``you stayed upright and moved forward, good'' or ``you fell over, bad''.

\emph{Reinforcement learning} (RL) learns from exactly this kind of feedback. The robot tries something, observes what happens, receives a number telling it how good that was, and slowly changes its behaviour so that the numbers get higher. This is how a child learns to ride a bicycle, and it is how many modern humanoid robots learn to walk, get up and kick.

\section*{The reinforcement learning loop}
Every RL problem has the same structure, shown in Figure~\ref{fig:rl1_loop}. There are two players.

\begin{itemize}[nosep]
  \item The \textbf{agent} is the learner and decision maker. For a robot this is its controller, often a neural network.
  \item The \textbf{environment} is everything else: the field, the ball, the other players and the physics of the world (or a simulator that imitates it).
\end{itemize}

\begin{center}
\begin{tikzpicture}[node distance=5.2cm, box/.style={draw,rounded corners,thick,minimum width=3.3cm,minimum height=1.6cm,align=center,font=\small}]
  \node[box,fill=blue!8] (ag) {\textbf{Agent}\\[-1pt]{\footnotesize the robot's controller}};
  \node[box,fill=green!8,right=of ag] (env) {\textbf{Environment}\\[-1pt]{\footnotesize field, ball, physics}};
  \draw[-{Stealth},thick] ([yshift=0.5cm]ag.east) -- node[above,font=\small]{action $a_t$} ([yshift=0.5cm]env.west);
  \draw[-{Stealth},thick,red!70!black] ([yshift=-0.05cm]env.west) -- node[above,font=\small,pos=0.5]{reward $r_{t+1}$} ([yshift=-0.05cm]ag.east);
  \draw[-{Stealth},thick] ([yshift=-0.55cm]env.west) -- node[below,font=\small]{new state $s_{t+1}$} ([yshift=-0.55cm]ag.east);
\end{tikzpicture}
\captionof{figure}{The agent--environment loop. In every time step the agent looks at the state, chooses an action, and the environment answers with a new state and a reward.}\label{fig:rl1_loop}
\end{center}

In every time step $t$ the following happens:
\begin{enumerate}[nosep]
  \item The agent observes the \textbf{state} $s_t$ (for example the position of the ball and the angles of its joints).
  \item The agent chooses an \textbf{action} $a_t$ (for example a target angle for each joint).
  \item The environment reacts. It returns a \textbf{reward} $r_{t+1}$ (a single number, positive for good, negative for bad) and the next state $s_{t+1}$.
\end{enumerate}
The agent's rule for choosing actions is called its \textbf{policy} $\pi$. A sequence of time steps from a start until the end (the goal is reached or the robot falls) is an \textbf{episode}. Learning means improving the policy.

\begin{keyidea}
Reinforcement learning = \textbf{state} $\to$ \textbf{action} $\to$ \textbf{reward}, repeated many times. The agent is never told the right action. It only gets rewards and has to discover for itself which actions lead to high rewards.
\end{keyidea}

\begin{center}
\renewcommand{\arraystretch}{1.3}
\begin{tabular}{p{2.6cm}p{3.6cm}p{3.4cm}p{3.5cm}}
\textbf{Task} & \textbf{State} & \textbf{Action} & \textbf{Reward}\\ \hline
Walking & joint angles and speeds, body tilt, wanted walking speed & target angle of each leg joint & $+$ for moving at the wanted speed, $-$ for falling\\
Penalty kick & ball position, robot pose & kick direction and strength & $+1$ for a goal, $0$ otherwise\\
Goalie & ball position and velocity & dive left, right, or stay & $+1$ for a save, $-1$ for a goal\\
\end{tabular}
\end{center}

\begin{example}
A goalie robot sees the ball coming. The \emph{state} is the position and velocity of the ball. The \emph{actions} are ``dive left'', ``stay'' and ``dive right''. If the ball is stopped the reward is $+1$, if it passes the robot the reward is $-1$. After many shots the robot learns, for example, that for a fast ball at the left it must dive left \emph{immediately}.
\end{example}

\section*{Return and discounting}
A single reward is not what the agent tries to maximise. It wants a high reward \emph{over the whole episode}. The sum of the rewards from time $t$ on is called the \textbf{return}. Rewards that come far in the future are usually worth less than rewards now, so we multiply the reward after $k$ further steps by $\gamma^k$, where $0<\gamma\le1$ is the \textbf{discount factor}:
\[
G_t = r_{t+1}+\gamma\,r_{t+2}+\gamma^2\,r_{t+3}+\dots=\sum_{k=0}^{\infty}\gamma^k\,r_{t+k+1}.
\]
With $\gamma$ close to $1$ the agent is patient, with $\gamma$ close to $0$ it only cares about the next reward.

\begin{center}
\begin{tikzpicture}
\begin{axis}[gr, axis lines=left, xlabel near ticks, ylabel near ticks, width=9cm,height=5.6cm, xlabel={steps $k$ into the future}, ylabel={weight $\gamma^k$}, xmin=0,xmax=20,ymin=0,ymax=1.05, xtick={0,5,10,15,20}, ytick={0,0.25,0.5,0.75,1}, legend pos=north east]
  \addplot[blue,domain=0:20,samples=21,mark=*,mark size=1.2pt]{0.5^x};
  \addlegendentry{$\gamma=0.5$}
  \addplot[red,domain=0:20,samples=21,mark=*,mark size=1.2pt]{0.9^x};
  \addlegendentry{$\gamma=0.9$}
  \addplot[green!50!black,domain=0:20,samples=21,mark=*,mark size=1.2pt]{0.99^x};
  \addlegendentry{$\gamma=0.99$}
\end{axis}
\end{tikzpicture}
\captionof{figure}{How much a reward $k$ steps in the future counts. With $\gamma=0.5$ it is almost worthless after $5$ steps, with $\gamma=0.99$ it still counts $95\,\%$ after $5$ steps.}\label{fig:rl1_gamma}
\end{center}

\begin{example}
The robot receives the rewards $-1,-1,-1,+10$ in four steps. With $\gamma=0.9$
\[
G=-1+0.9\cdot(-1)+0.9^2\cdot(-1)+0.9^3\cdot 10=-1-0.9-0.81+7.29=4.58 .
\]
With $\gamma=1$ (no discount) the return is $-1-1-1+10=7$, with $\gamma=0.5$ it is only $-1-0.5-0.25+1.25=-0.5$: an impatient agent sees the three costly steps but hardly values the far reward.
\end{example}

\begin{aibox}
Training a walking policy typically uses $\gamma$ close to $1$ (for example $0.99$), because the reward for staying upright must be linked to the actions that happened many steps earlier. The same idea, ``the future counts a bit less'', appears in finance as interest.
\end{aibox}

\section*{A tiny world: the corridor}
We will use one very small world throughout this part. A robot stands in a corridor of five cells. Cell~5 holds the ball (the goal). The robot can step \textbf{left} (L) or \textbf{right} (R). At the left wall it just stays where it is. Every step costs $-1$ (it needs time and energy), and the step that reaches the goal gives $+10$. The episode ends when the goal is reached.

\begin{center}
\begin{tikzpicture}[cell/.style={draw,thick,minimum size=1.5cm,font=\large}]
  \foreach \i in {1,...,5}{
    \node[cell] (c\i) at (1.6*\i,0) {};
    \node[font=\footnotesize,gray] at (1.6*\i,-1.05) {cell \i};
  }
  \fill[green!25] (c5.south west) rectangle (c5.north east);
  \draw[thick] (c5.south west) rectangle (c5.north east);
  \node[font=\Large] at (c5) {$\bigcirc$};
  \node[font=\Large,blue!70!black] at (c1) {\textbf{R}};
  \draw[-{Stealth},thick,blue!70!black] (c1.north east) to[bend left=35] node[above,font=\scriptsize]{R: $-1$} (c2.north west);
  \draw[-{Stealth},thick,blue!70!black] (c4.north east) to[bend left=35] node[above,font=\scriptsize]{R: $+10$} (c5.north west);
\end{tikzpicture}
\captionof{figure}{The corridor world. The robot (letter R) starts in cell~1 and must reach the goal (circle) in cell~5. Each step costs $-1$, the last step gives $+10$.}\label{fig:rl1_corridor}
\end{center}

\begin{example}
The robot walks right four times. The rewards are $-1,-1,-1,+10$ and, with $\gamma=0.9$, the return is $4.58$ (see above). Suppose instead it first steps left (it stays in cell~1 and pays $-1$) and then walks right four times. The rewards are $-1,-1,-1,-1,+10$:
\[
G=-1-0.9-0.81-0.729+0.9^4\cdot 10=-3.439+6.561=3.122 .
\]
The wasted step lowers the return from $4.58$ to $3.122$. This is how RL punishes detours without ever telling the robot ``walking right is correct''.
\end{example}

\section*{Policies}
A \textbf{deterministic policy} says exactly which action to take in each state, written $a=\pi(s)$. A \textbf{stochastic policy} gives a probability for each action, written $\pi(a\mid s)$. For the corridor, ``always go right'' is
\[
\pi(\text{cell }1)=\pi(\text{cell }2)=\pi(\text{cell }3)=\pi(\text{cell }4)=R .
\]
A stochastic policy could be ``go right with probability $0.8$, left with probability $0.2$''. Randomness is useful during learning, because the robot has to \emph{try} things.

\section*{Exploration versus exploitation}
Imagine a striker with three different kicks, A, B and C. It does not know which one scores most often. Each kick is a \textbf{slot machine} (a ``bandit''): the reward is $1$ for a goal and $0$ otherwise. The robot keeps an estimate $Q$ of the success rate of every kick. After the $n$-th try of a kick with reward $r$ the new average is
\[
Q_{n}=Q_{n-1}+\frac{1}{n}\bigl(r-Q_{n-1}\bigr).
\]
This is the mean of all rewards so far, written as ``old value plus a correction towards the new reward''. The same pattern will return in Q-learning in Chapter~\ref{ch:rl3}.

The robot faces a dilemma. It can \textbf{exploit}: always use the kick with the best estimate. Or it can \textbf{explore}: try other kicks to learn whether they are better. The simplest compromise is the \textbf{$\varepsilon$-greedy} rule: with probability $1-\varepsilon$ take the best-known action, with probability $\varepsilon$ a random one.

\begin{example}
The robot tries six kicks and observes (kick, reward): $(A,1)$, $(B,0)$, $(C,1)$, $(C,1)$, $(A,0)$, $(C,0)$. The estimates after each try are:
\begin{center}
\renewcommand{\arraystretch}{1.2}
\begin{tabular}{c|ccc}
after try & $Q_A$ & $Q_B$ & $Q_C$\\ \hline
1 & $1$ & $0$ & $0$\\
2 & $1$ & $0$ & $0$\\
3 & $1$ & $0$ & $1$\\
4 & $1$ & $0$ & $1$\\
5 & $0.5$ & $0$ & $1$\\
6 & $0.5$ & $0$ & $0.667$
\end{tabular}
\end{center}
(Kick B was tried once with reward $0$, so $Q_B=0$.) The greedy choice is now kick C. But B was tried only once. Its success rate may actually be high, and only exploration can find out.
\end{example}

\begin{keyidea}[Exploration versus exploitation]
Too much exploitation: the robot sticks with the first thing that worked and may never find the best behaviour. Too much exploration: it wastes time on bad actions. Almost every RL method uses a rule such as $\varepsilon$-greedy and makes $\varepsilon$ smaller during training.
\end{keyidea}

\section*{Practice problems}
\begin{enumerate}[label=\textbf{\thechapter.\arabic*}]
  \item For each situation name the agent, the state, the action and a sensible reward: (a) a robot learning to get up from the ground, (b) a robot dribbling the ball towards the goal, (c) a thermostat that heats a room.
  \item Compute the return for the rewards $2,2,2,2$ with $\gamma=0.9$.
  \item The rewards are $0,0,5$ and $\gamma=0.5$. Compute the return $G_t$.
  \item A reward of $10$ after $3$ steps shall be worth exactly as much as a reward of $8$ right now. Which $\gamma$ is needed?
  \item In the corridor of Figure~\ref{fig:rl1_corridor} the robot starts in cell~1 and always walks left. The episode is stopped after $10$ steps. Compute its return for $\gamma=0.9$.
  \item Explain why a reward of $+1$ only when a goal is scored makes it hard for a robot to learn to walk towards the ball. Suggest one additional reward that would help.
  \item The three-kick striker tries kick $C$ five times and gets $0,1,1,0,1$. Compute $Q_C$ after each try with the update $Q_n=Q_{n-1}+\frac1n(r-Q_{n-1})$.
  \item An $\varepsilon$-greedy robot has three actions and $\varepsilon=0.1$ (a random action is chosen uniformly among all three, including the greedy one). With which probability is the greedy action played? With which probability each of the other two?
  \item \texttt{[code]} Simulate the three kicks with true success rates $0.3$, $0.5$, $0.8$. Use an $\varepsilon$-greedy rule with $\varepsilon=0.1$ for $1000$ tries and print the final estimates $Q_A,Q_B,Q_C$. Which kick does the robot use most often?
  \item[\textbf{\thechapter.10}$\star$] A reward of $1$ is received in every step forever. Show that the return is $\frac{1}{1-\gamma}$ for $0<\gamma<1$ and compute it for $\gamma=0.99$. (Hint: geometric series.)
\end{enumerate}
